\documentclass[main.tex]{subfiles}
\begin{document}
	

\section{Normal Covers}\label{sec:ch2_normal_covers}

{\definition{} 
A function $\rho : X \times X \to [0, \infty)$ is a 
\uwave{pseudo-metric}\index{pseudo-metric| \S\ref{sec:ch2_normal_covers}\hfill}
 on $X$ if for all $x, y, z \in X$, we have $\rho(x, x) = 0$, $\rho(x, y) = \rho(y, x)$, 
 and $\rho(x, z) \le \rho(x, y) + \rho(y, z)$. A pseudo-metric $\rho$ is a 
\uwave{metric}\index{metric| \S\ref{sec:ch2_normal_covers}\hfill} 
 if $\rho(x, y) = 0$ implies $x = y$. 
 The topology $\tau_\rho$ on $X$ generated by a pseudo-metric $\rho$ is 
 the topology which has the family $\mathcal{B} = \{B_\rho(x, r) : x \in X, r > 0\}$ 
 of balls as a base, where $B_\rho(x, r) = \{y \in X : \rho(x, y) < r\}$.
}

Tukey [1940]\cite{Tukey} proved that every metric space is fully normal.

{\theorem. Every pseudo-metric space is fully normal.
}

\begin{proof}
Let $(X, \rho)$ be a pseudo-metric space. Let $\mathcal{U}$ be an open cover of $X$. For each $x \in X$, pick some $U_x \in \mathcal{U}$ and 
an integer $m(x)\ge 1$ such that $x \in U_x$ and

\begin{equation}
	\mathclap{B_\rho\left(x, \frac{3}{m(x)}\right) \subset U_x\in \mathcal{U}.} \tag*{(1)}
\end{equation}

Let $\mathcal{V}$ be the sets of smaller balls $B_\rho(x, \frac{1}{m(x)})$
for all  $x \in X$ and its chosen $m(x)$ satisfying Condition (1).
Then $\mathcal{V}$ is an open refinement of $\mathcal{U}$. 
We will show that it is also a point-star refinement of $\mathcal{U}$, 
so that $X$ is fully normal. 
Let $x \in X$, then there exits at least one $u \in X$ such that
$x\in B_\rho\left(u, \frac{1}{m(u)}\right)\in\mathcal{V}$
since $\mathcal{V}$ covers $X$.
So, we are able to define
$$k(x) = \min\{m(u) : u\in X, x\in B_\rho\left(u, \frac{1}{m(u)}\right)\}.$$
and pick $u_x\in X$ such that $k(x) = m(u_x)$. Now, we have

\begin{equation}
	\mathclap{x\in B_\rho\left(u_x, \frac{1}{m(u_x)}\right)\in \mathcal{V}.} \tag*{(2)}
\end{equation}

It suffices to show that 

\begin{equation}
	\mathclap{\mathrm{st}(x, \mathcal{V}) \subset U_{u_x}.} \tag*{(3)}
\end{equation}

By (1), it suffices to show
\begin{equation}
	\mathclap{\mathrm{st}(x, \mathcal{V}) \subset B_\rho(u_x, \frac{3}{m(u_x)}).} \tag*{(4)}
\end{equation}

For if $y \in \mathrm{st}(x, \mathcal{V})$, then exits $t\in X$ 
such that $y, x \in B_\rho\left(t, \frac{1}{m(t)}\right)$.
Since $m(u_x)=k(x) \le m(t)$, we have

$$\rho(y, u_x) \le \rho(y, t) + \rho(t, x) + \rho(x, u_x) < \frac{3}{m(u_x)}.$$
Thus (4) holds. 
\end{proof}

{\theorem{} \label{th:ch2_1_3_dowker}
\textup{(Dowker-Morita)} A countable functionally open cover of $X$ 
has a countable star-finite functionally open refinement.
}

\begin{proof} 
Let $\mathcal{U} = \{U_n : n \ge 1\}$ be a functionally open cover of $X$, 
where each $U_n = f_n^{-1}(0, 1]$ for some continuous function $f_n : X \to [0, 1]$. 
Define a continous function $f : X \to [0, 1]$ by
$$f(x) = \sum_{n=1}^{\infty} f_n(x) / 2^n, x\in X.$$
For $n \ge 1$, let $G_n = f^{-1}\left(\frac{1}{2^n}, 1\right]$, 
$F_n = f^{-1}\left[\frac{1}{2^n}, 1\right]$, and set $F_0 = \emptyset$.
It is easy to check that $G_{n+1} \subset \bigcup_{k=1}^{n} U_k$ for each $n \ge 1$.
Let
$$\mathcal{W} = \{W_{n,k}: W_{n,k} = U_k \bigcap (G_{n+1} \setminus F_{n-1})
 \text{ for } 
1\le k \le n<\omega\}.$$
To see that $\mathcal{W}$ covers $X$, let $x \in X$. 
Since $\{G_n\}$ covers $X$, $f(x) > 0$. Let $m$ be the smallest integer such that 
$x \in F_m$. Since $F_m \subset G_{m+1} \subset \bigcup_{k \le m} U_k$, 
there exists $k \le m$ such that 
$x \in U_k \bigcap (G_{m+1} \setminus F_{m-1}) = W_{m,k}$.
Hence $\mathcal{W}$ is a countable functionally open refinement of $\mathcal{U}$.
To see that $\mathcal{W}$ is star-finite, fix $(k, j)$, $1 \le j \le k$. 
For any $(m. i)$, $1 \le i \le m$, if $m > k + 2$, 
then $W_{k,j} \subset G_{k+1} \subset F_{k+1} \subset F_{m-1}$. 
Then $W_{k,j} \bigcap W_{m,i} \subset F_{m-1}\bigcap(X\setminus F_{m-1})= \emptyset$.
Hence if $W_{m,j} \bigcap W_{k,i} \neq \emptyset$, then $m < k+2$.
It follows that $(\mathcal{W})_{W_{k,j}}\subset \{W_{m,i} : m < k+2, 1 \le i \le m\}$, 
so $\mathcal{W}$ is star-finite.
\end{proof}

{\theorem{\label{th:ch2_1_Frink}} \textup{(Frink [1937]\cite{Frink})}. Suppose $f: X \times X \to [0, \infty)$ which satisfies

\textup{(a)} $f(x,x)=0$, and $f(x,y)=f(y,x)$,

\textup{(b)} For any $r > 0$ and any $x, y\in X$,  we have $f(x,y) < 2r$ if $f(x,z) < r$ 
and $f(z, y) < r$ for any $z\in X$.

Then there is a pseudo-metric $\rho$ on $X$ such that for any $x, y \in X$,

\textup{(i)}  $\rho(x,y) \le f(x,y)$,

\textup{(ii)}  $f(x,y) \le 4\rho(x,y)$.
}

\begin{proof} 
For $x, y\in X$, let
\begin{gather}
	g_0(x,y,\emptyset) = f(x,y), \\
	g_1(x,y,s) = f(x,a_1) + f(a_1,y), \text{ for } s = (a_1) \in X^1\\
	g_n(x,y,s) = f(x,a_1) + \sum_{i=1}^{n-1} f(a_i, a_{i+1}) + f(a_n,y), \text{ for } 
	   n \ge 2, s = (a_1, \ldots, a_n) \in X^n, \\
	\rho(x,y) = \inf \{ g_n(x,y,s) : s \in X^n, n \ge 0 \}.
\end{gather}
It is easily verified that $\rho$ is a pseudo-metric on $X$ and $\rho(x,y) \le f(x,y)$. It remains to show (ii). For $x, y \in X$, let 
\begin{gather}
	h_0(x,y,\emptyset) = f(x,y), \\
    h_1(x,y,s) = 2f(x,a_1) +2f(a_1,y), \text{ for } s = (a_1) \in X^1\\
	h_n(x,y,s) = 2f(x,a_1) + 4\sum_{i=1}^{n-1} f(a_i, a_{i+1}) + 2f(a_n,y), \text{ for } 
	n \ge 2, s = (a_1, \ldots, a_n) \in X^n.
\end{gather}

\textsc{Claim.}  As usual, we assume that $X^0 = \{\emptyset\}$. Then
\begin{gather}
f(x,y) \le h_n(x,y,s), \textup{ for any } x, y \in X, n \in\omega \textup{ and } s \in X^n.
\label{gather:ch2_1_claim}
\end{gather}

If we have proved Claim \eqref{gather:ch2_1_claim}, then $f(x,y) \le 4g_n(x,y,s)$, so $f(x,y) \le 4\rho(x,y)$, that completes our proof.

For any $x, y, a \in X$ we have 
\begin{gather}
	f(x, y) \le 2 f(a, y) \quad \textup{or}\quad f(x, y) \le 2 f(x, a)  \label{eq:2.1.ineq}
\end{gather}

The proof of \eqref{eq:2.1.ineq} has two cases.

Case 1. If $f(x, a) \le f(a, y)$, then for any $r > f(a, y)$, $f(x, a) \le f(a, y) <  r$. 
By hypothesis (b), $f(x, y) < 2 r$. So, $f(x, y) \le 2 f(a, y)$ as  $r \rightarrow f(a, y)$ above.

Case 2. If $f(a, y) < f(x, a)$, then by (b) we have $f(x, y) < 2 r$ for any $r > f(x, a)$.
Again, it results that $f(x, y) \le f(x, a)$ as  $r \rightarrow f(x, a)$ above.

\textsc{Proof of Claim \eqref{gather:ch2_1_claim}}. We use induction. For $n=0$, we have $f(x, y) = h_0(x, y, \phi)$. 
Let $m \ge 1$, and assume that for all $i < m$, 
$s\in X^i$, $f(x, y) \le h_i(x, y, s)$. Let us show Claim \eqref{gather:ch2_1_claim}
for $n\ge1$.

%(2) $f(x, y) \le h_{m}(x, y, s)$ for $x, y \in X$, $s = (a_{1}, \dots, a_{m})\in X^m$.

In Case 1 of \eqref{eq:2.1.ineq}, 
$f(x, y) \le 2 f(a_{m}, y) \le 2 g_{m}(x, y, s) \le h_{m}(x, y, s)$, we are finished.

Now, we turn to Case 2 of  \eqref{eq:2.1.ineq}: $f(x, y) \le 2 f(x, a_{m})$. 
Let 
$$k = \min \{i \le m : f(x, y) \le 2 f(x, a_{i})\}.$$
Then $1 \le k \le m$, $f(x, y) \le 2 f(x, a_{k})$. If $k = 1$, then 
$$f(x, y) \le 2 f(x, a_{1}) \le 2 g_{m}(x, y, s) \le h_{m}(x, y, s),$$
we are finished. 
So we assume $k \ge 2$. By the definition of $k$ and \eqref{eq:2.1.ineq}, 
$f(x, y) \le 2 f(a_{k-1}, y)$. It follows that

$$f(x, y) \le f(x, a_{k}) + f(a_{k-1}, y).$$
Let $s_{1} = (a_{1}, \dots, a_{k-1})$, $s_{2} = (a_{k}, \dots, a_{m})$, and $s = (a_{1}, \dots, a_{k}, \dots, a_{m})$. Since $k - 1 < m$, $m - k + 1 < m$, by the induction hypothesis, we have
$$f(x, a_{k}) \le f_{k-1}(x, a_{k}, s_{1}), \quad
f(a_{k-1}, y) \le f_{m-k+1}(a_{k-1}, y, s_{2}).$$
Then 
\begin{equation*}
	\begin{aligned}
f(x, y)    &\le f_{k-1}(x, a_{k}, s_{1}) + f_{m-k+1}(a_{k-1}, y, s_{2}) \\
              &= 2 f(x, a_{1}) + 4 \sum_{i=1}^{k-2} f(a_{i}, a_{i+1}) + 2 f(a_{m}, y) \\
	          &= h_{m}(x, y, s)
	 \end{aligned}
\end{equation*}
This completes the induction and the Claim.
\end{proof}


{\theorem{}\label{th:ch2_1_5_tukey}
\textup{(Tukey [1940])\cite{Tukey}}. Suppose there exists 
a sequence $\{\mathcal{U}_n: n\in\omega\}$ of open covers of a space $(X, \tau)$ 
such that $\mathcal{U}_0 = \{X\}$ and 
$\mathcal{U}_{n+1}$ is a point star refinement of $\mathcal{U}_n$ for each $n \ge 1$. 
Then there exists a pseudo-metric $\rho$ on $X$ such that for each $x \in X, n \in \omega:$

\textup{(i)}  $\operatorname{st}(x, \mathcal{U}_n) \subset B_\rho(x, 1/{2^n})$,
 and $\tau_\rho \subset \tau$ where  $\tau_\rho$ is the topology generated by metric $\rho$.

\textup{(ii)} $B_\rho(x, 1/{2^{n+2}}) \subset \operatorname{st}(x, \mathcal{U}_n)$
}

\begin{proof}
Let $\{\mathcal{U}_n: n\in\omega\}$ be as in the hypothesis. 
Define $f: X \times X \to [0, \infty)$ by:

$$f(x, y) =  \begin{cases} 
	 %1, & \text{if } y \notin \bigcup_{n < \omega} \operatorname{st}(x, \mathcal{U}_n), \\  
	 0, & \text{if } y \in \bigcap_{n < \omega} \operatorname{st}(x, \mathcal{U}_n), \\  
	 \frac{1}{2^n}, &
	   \text{if } y \in \operatorname{st}(x, \mathcal{U}_{n-1}) 
	   \setminus \operatorname{st}(x, \mathcal{U}_{n}), n \ge 1.  
\end{cases}$$

Clearly, $f$ satisfies the condition (a) of Theorem \ref{th:ch2_1_Frink}.

To show the condition (b), first, it is easy to see that
for $x, y \in X, n \in \omega$
\begin{gather}
	y \in \operatorname{st}(x, \mathcal{U}_n) \iff f(x, y) < \frac{1}{2^n}. \label{eq:ch2_1_tukey}
\end{gather}
Use the notation in condition (b) of Theorem \ref{th:ch2_1_Frink} to
assume that $x, y, z \in X, r > 0$ such that $f(x, z) < r$ and $f(z, y) < r$.
%Suppose $x, y, z \in X, r > 0$ such that $f(x, y) < r$ and $f(y, z) < r$. 
We show $f(x, y) < 2r$ in
two cases.

Case 1. Suppose $r \ge 1/2$. If $y \in \operatorname{st}(x, \mathcal{U}_n)$ 
for each $n \in \omega$, then $f(x, y) = 0 < 2r$. Otherwise there is a least $k \ge 1$
 such that $y \in \operatorname{st}(x, \mathcal{U}_{k-1}) \setminus \operatorname{st}(x, \mathcal{U}_k)$. Then $f(x, y) = 1/{2^k} \le 1/2 < 2r$, contradiction.

Case 2. Suppose $r <1/2$. Let $k = \min\{n \in \omega : 1/{2^n} < r\}$.
Then $1/{2^k} < r < 1/2$, so $k > 1$. Either $f(x, y) = 0 < 1/{2^{k-1}}$ or 
$f(x, z) = 1/{2^m} < r$, so $k \le m$. Hence $f(x, z) \le 1/{2^m} < 1/2^{k-1}$.
For the same reason we also have $f(z, y) < 1/{2^{k-1}}$. 
It follows from \eqref{eq:ch2_1_tukey} that $z \in \operatorname{st}(x, \mathcal{U}_{k-1})$ 
and $y \in \operatorname{st}(z, \mathcal{U}_{k-1})$. 
There are $U_1, U_2 \in \mathcal{U}_{k-1}$ such that $x, z \in U_1$ 
and $z, y \in U_2$. Since $z \in U_1 \cap U_2$ 
and $\mathcal{U}_{k-1}$ is a point star refinement of $\mathcal{U}_{k-2}$, 
$U_1 \cup U_2 \subset U$ for some $U \in \mathcal{U}_{k-2}$. 
Then $y \in \operatorname{st}(x, \mathcal{U}_{k-2})$. 
By \eqref{eq:ch2_1_tukey}, 
$f(x, y) < 1/{2^{k-2}} < 2r$.
Hence condition (b) holds of Theorem \ref{th:ch2_1_Frink} holds.

By Theorem \ref{th:ch2_1_Frink}, there exists a pseudo-metric $\rho$ on $X$ such that for each $x, y \in X$,
 \begin{gather}
 \rho(x, y) \le f(x, y), \quad
f(x, y) \le 4\rho(x, y). \label{eq:th_frink_f_rho}
\end{gather}

It follows from \eqref{eq:th_frink_f_rho} and \eqref{eq:ch2_1_tukey} that:

\textup{(i)} $\operatorname{st}(x, \mathcal{U}_n) \subset B_\rho(x,1/{2^n})$ and $\tau_\rho \subset \tau$.

\textup{(ii)} $B_\rho(x, 1/{2^{n+2}}) \subset \operatorname{st}(x, \mathcal{U}_n)$.
\end{proof}

{\theorem{} \label{th:ch2_1_6_tukey}
	\textup{(Tukey [1940]\cite{Tukey}, 
		Morita [1962, 1964]\cite{Morita1962}\cite{Morita1964})}. 
	For any open cover $\mathcal{U}$ of a space $(X, \tau)$, 
	the following are equivalent:

\textup{(i)}   $\mathcal{U}$ is a normal cover.

\textup{(ii)}  $\mathcal{U}$ has a locally finite  functionally open refinement 
which is also $\sigma$-discrete.

\textup{(iii)}  $\mathcal{U}$ has an locally finite functionally open refinement.

\textup{(iv)}  $\mathcal{U}$ has a $\sigma$-locally finite functionally open refinement.

\textup{(v)}  There is a partition of unity $\{\varphi_\alpha : \alpha \in A\}$ on $X$ 
such that $\{\varphi_\alpha^{-1}((0, 1]) : \alpha \in \Delta\}$ refines $\mathcal{U}$.
}

\begin{proof} (i) $\Rightarrow$ (ii).
Suppose $\mathcal{U}$ is normal and $\{\mathcal{U}_n: n\in\omega\}$ is a sequence of 
open covers of $(X, \tau)$ such that $\mathcal{U}_0 = \{X\}$, 
$\mathcal{U}_1$ 
refines $\mathcal{U}$ and $\mathcal{U}_{n+1}$ point-star refines 
$\mathcal{U}_n$ for each $n \ge 1$.
By Theorem \ref{th:ch2_1_5_tukey} there exists a pseudo-metric $d$ on $X$ such that:

1. $\tau_d \subset \tau$.

2. $B_\rho(x, 1/{2^{n+2}}) \subset \text{st}(x, \mathcal{U}_n), \quad x \in X, n < \omega$.

\noindent Since $\{\operatorname{st}(x, \mathcal{U}_2) : x \in X\}$ refines $\mathcal{U}_1$, 
$\mathcal{B} = \{B_d(x, 1/{16}) : x \in X\}$ is a refinement of $\mathcal{U}$. 
By Stone's Theorem (see Theorem 4.4.1 in Engelking [1977]\cite{Engelking}),  % TO_DO
$\mathcal{B}$ has a locally finite  and $\sigma$-discrete open refinement 
$\mathcal{V}$ in $(X, \tau_d)$. Since pseudo-metric spaces are perfectly normal, 
every member of $\mathcal{V}$ is functionally open. 
Since $\tau_\rho \subset \tau$, $\mathcal{V}$ is a locally finite 
and $\sigma$-discrete functionally open refinement of $\mathcal{U}$ in $(X, \tau)$.

(ii) $\Rightarrow$ (iii) $\Rightarrow$ (iv)
Clear.

(iv) $\Rightarrow$ (iii). 
Assume (iv). Let $\mathcal{U} = \{U_\alpha: \alpha \in \Delta\}$ be an open cover of $X$. 
Let $\mathcal{V} = \bigcup_{n < \omega} \mathcal{V}_n$ be 
a refinement of $\mathcal{U}$ where each $\mathcal{V}_n = \{V_{n,\alpha}: \alpha \in \Delta\}$
 is locally finite functionally open with $V_{n,\alpha} \subset U_\alpha$.


By Lemma \ref{lemma:prelim_3_8} and Theorem \ref{th:ch2_1_3_dowker} 
there is a countable star-finite functionally open cover 
$\mathcal{W}_0 = \{W_n\}_{n < \omega}$ such that $\mathcal{W}_0$ 
refines $\mathcal{V}_n$ for each $n$.

Let


$$\mathcal{H} = \{W_n \cap V_{n,\alpha} : n < \omega, \alpha \in \Delta\}.$$

It is easy to see that $\mathcal{H}$ is a LF functionally open refinement of $\mathcal{U}$.

**(iii) $\rightarrow$ (v)**

Suppose $X$ has a LF functionally open ...

Here is the extracted and cleaned transcription of the mathematical text and proof shown in **jiang26.pdf**:
\end{proof} 
---

refinement $\mathcal{V} = \{V_\alpha\}_{\alpha \in \Delta}$ of $\mathcal{U}$. By Lemma 1.3.8 there is a partition of unity $\{g_\alpha\}_{\alpha \in \Delta}$ such that $g_\alpha^{-1}((0, 1]) \subset V_\alpha \subset U_\alpha$ for each $\alpha \in \Delta$. Hence $\{g_\alpha^{-1}((0, 1]) : \alpha \in \Delta\}$ refines $\mathcal{U}$.

**(v) $\rightarrow$ (i)**

Let $\{g_\alpha\}_{\alpha \in \Delta}$ be a partition of unity on $(X, \tau)$ such that $\{g_\alpha^{-1}((0, 1]) : \alpha \in \Delta\}$ refines $\mathcal{U}$, where each $\varphi_\alpha : X \rightarrow [0, 1]$ is continuous with $\sum_{\alpha \in \Delta} \varphi_\alpha(x) = 1$. Let

$$Y = \{u \in I^\Delta : \{\alpha \in \Delta : u(\alpha) > 0\} \text{ is finite}\}.$$

Let $\pi_\alpha : I^\Delta \rightarrow I_\alpha = I$ be projections: $\pi_\alpha(u) = u(\alpha)$. Define a function $d : Y \times Y \rightarrow [0, \infty)$ by

$$d(u, v) = \sum_{\alpha \in \Delta} \vert{}u - v\vert{}(\alpha) = \sum_{\alpha \in \Delta} \vert{}\pi_\alpha(u) - \pi_\alpha(v)\vert{}.$$

It is easy to see that $d$ is a metric on $Y$. Define $f : X \rightarrow Y$ by $\pi_\alpha \circ f = g_\alpha$ for $\alpha \in \Delta$. Let $M = f(X)$. Then $M \subset Y$.

*Claim.* $f : (X, \tau) \rightarrow (M, \tau_d\vert{}_M)$ is continuous, where $\tau_d$ is the topology generated by metric $d$.

*Proof.* Let $x_0 \in X$ and $\epsilon > 0$. Let $\Gamma = \{\alpha \in \Delta : \varphi_\alpha(x_0) > 0\}$. We have:

1. $\sum_{\alpha \notin \Gamma} \varphi_\alpha(x_0) = 0$.



Suppose $\Gamma = \{\alpha_1, \dots, \alpha_m\}$. For each $i \le m$, there is an open neighborhood $V_i$ of $x_0$ such that if $x \in V_i$, then $\vert{}\varphi_{\alpha_i}(x) - \varphi_{\alpha_i}(x_0)\vert{} < \epsilon / 2m$. Let $V = \bigcap_{i=1}^m V_i$. Then for each $x \in V$:
2. $\sum_{\alpha \in \Gamma} \vert{}\varphi_\alpha(x) - \varphi_\alpha(x_0)\vert{} < \epsilon / 2$.

By (1), we have:


$$\sum_{\alpha \notin \Gamma} \varphi_\alpha(x) + \sum_{\alpha \in \Gamma} \varphi_\alpha(x_0) = \sum_{\alpha \in \Delta} \varphi_\alpha(x) = 1 = \sum_{\alpha \in \Gamma} \varphi_\alpha(x_0).$$

By (2), we have:
3. $\sum_{\alpha \notin \Gamma} \varphi_\alpha(x) = \sum_{\alpha \in \Gamma} (\varphi_\alpha(x_0) - \varphi_\alpha(x)) < \epsilon / 2$.

Then we have for each $x \in V$ ...

Here is the extracted and cleaned transcription of the mathematical text and proof shown in **jiang27.pdf**:

---

$$d(f(x), f(x_0)) = \sum_{\alpha \in \Delta} \vert{}f(x)(\alpha) - f(x_0)(\alpha)\vert{}$$

$$= \sum_{\alpha \in \Gamma} \vert{}\varphi_\alpha(x) - \varphi_\alpha(x_0)\vert{} + \sum_{\alpha \notin \Gamma} \vert{}\varphi_\alpha(x) - \varphi_\alpha(x_0)\vert{} < \epsilon.$$

This completes the proof of the claim.

For $\alpha \in \Delta$, let $H_\alpha = \{u \in M : u(\alpha) > 0\}$. Then:
(4) For each $\alpha \in \Delta$, $f^{-1}(H_\alpha) = g_\alpha^{-1}((0, 1]) \subset U_\alpha$.
(5) $\mathcal{H} = \{H_\alpha : \alpha \in \Delta\}$ is an open cover of $(M, \tau_d\vert{}_M)$.

*Proof.* For each $x \in X$, pick $\alpha_0 \in \Delta$ such that $g_{\alpha_0}(x) > 0$. Then $(\pi_{\alpha_0} \circ f)(x) \in (0, 1]$, so $f(x) \in H_{\alpha_0}$, and $\mathcal{H}$ covers $M$. For each $\alpha \in \Delta$, it is easy to see $H_\alpha \in \tau_d\vert{}_M$. Let $f(x) \in H_\alpha$, where $x \in X$. Let $r_x = (\pi_\alpha \circ f)(x) = \varphi_\alpha(x) \in (0, 1]$, so $r_x > 0$. It suffices to show:

(6) $B_d(f(x), r_x) \cap M \subset H_\alpha$.

*Proof.* Suppose $f(y) \in B_d(f(x), r_x) \cap M$, where $y \in X$. Then

$$r_x > d(f(y), f(x)) = \sum_{\beta \in \Delta} \vert{}\varphi_\beta(y) - \varphi_\beta(x)\vert{} \ge \vert{}\varphi_\alpha(y) - \varphi_\alpha(x)\vert{},$$


hence $\pi_\alpha(f(y)) = \varphi_\alpha(y) > 0$, so $f(y) \in H_\alpha$ and (6) holds.

By Theorem 2.1.4, $(M, \tau_d\vert{}_M)$ is fully normal. By (5), $\mathcal{H}$ is a normal cover of $M$. By the claim and (4), $f^{-1}(\mathcal{H}) = \{f^{-1}(H_\alpha) : \alpha \in \Delta\}$ is a normal cover of $(X, \tau)$, so $\mathcal{U}$ is a normal cover.

{\definition{} 
	\textup{(Bing [1951]\cite{Bing})} 
	$X$ is *
\uwave{$\lambda$-strongly screenable}
\index{$\lambda$-strongly screenable| \S\ref{sec:ch2_normal_covers}\hfill}
 $($\uwave{$\lambda$-screenable}
 \index{$\lambda$-screenable| \S\ref{sec:ch2_normal_covers}\hfill}$)$
if every open $\lambda$-cover of $X$ has a 
$\sigma$-discrete $(\sigma$-disjoint$)$ open refinement.
}

{\corollary{} 
	A fully normal space is $\sigma$-paracompact.
}


\end{document}